\documentclass[letter,11pt]{letter}

\usepackage[margin=0.89in]{geometry}
%\usepackage{moreverb,latexsym}
%\usepackage{graphics,graphicx,curves,epic,eepic,ulem,bbm,mathbbol,enumerate,latexsym,multicol}
\usepackage{lipsum}
\usepackage{amssymb}
\usepackage{amsmath}
\usepackage{amsfonts}
\usepackage{graphicx}
\usepackage{caption}
%\usepackage{epstopdf}
%\usepackage{theorem}
\usepackage{mathrsfs}
\usepackage{mathtools}
%\usepackage{empheq}
\usepackage{bm}
\usepackage{bbm}
\usepackage{color}
\usepackage{setspace}
\usepackage{exscale}
\usepackage{relsize}
%\usepackage{float}
\usepackage{picinpar}
\usepackage{extarrows}
\usepackage{multirow}
%\usepackage{cite}
\usepackage[backend=bibtex,style=numeric]{biblatex}
\addbibresource{ref.bib}
\defbibheading{bibliography}[\refname]{\par\bigskip\noindent\textbf{#1}\par}

\usepackage{nicefrac}
%%% PACKAGES
\usepackage{booktabs} % for much better looking tables
\usepackage{array} % for better arrays (eg matrices) in maths
\usepackage{paralist}
\usepackage[normalem]{ulem}
\newcommand{\vla}[1]{\textcolor{magenta}{#1}}
\usepackage{tikz}
\usetikzlibrary{positioning}
\usepackage{xcolor}
\usepackage{enumitem,kantlipsum}

%\usepackage{subcaption}
%\usepackage{epsfig}
\usepackage{algorithm}
\usepackage{algpseudocode}

\newcommand{\RIcolor}[1]{{\leavevmode\color{red} #1}}
\newcommand{\RIIcolor}[1]{{\leavevmode\color{blue} #1}}
\newcommand{\RIIIcolor}[1]{{\leavevmode\color{magenta} #1}}
\usepackage{amsthm}
\usepackage{graphicx}
\usepackage{hyperref}  % optional, for clickable references

% Define a custom figure counter for letters
\newcounter{myfigure}
\newcommand{\myfigure}[3][0.8\textwidth]{%
	\refstepcounter{myfigure}%
	\begin{center}
		\includegraphics[width=#1, height=0.75\textheight, keepaspectratio]{#2}%
		\\[0.5em]
		\textbf{Figure~\themyfigure:} #3
	\end{center}%
}


\theoremstyle{remark}
\newtheorem*{remark}{Remark}

\newcommand{\uvec}[2][3]{\bm{#2\mkern-#1mu}\mkern#1mu}




\newcommand*\xbar[1]{%
  \hbox{%
    \vbox{%
      \hrule height 0.5pt % The actual bar
      \kern0.4ex%         % Distance between bar and symbol
      \hbox{%
        \kern-0.05em%      % Shortening on the left side
        \ensuremath{#1}%
        \kern-0.05em%      % Shortening on the right side
      }%
    }%
  }%
}

% Some of the article customisations are relevant for this class

\name{Lorenzo Micalizzi} % To be used for the return address on the envelope
\signature{Lorenzo Micalizzi} % Goes after the closing (ie at the end of the letter, with space for a signature)
\address{Lorenzo Micalizzi\\
Department of Mathematics\\
North Carolina State University\\
North Carolina, USA}
% Alternatively, these may be set on an individual basis within each letter environment.

%\makelabels % this command prints envelope labels on the final page of the document

\begin{document}
\begin{letter}{Computer Methods in Applied Mechanics and Engineering\\
CMAME-D-26-02015}

\opening{Dear Editor:} % eg Hello.
We would like to thank the reviewers for carefully reading our paper and for their valuable suggestions. We have revised the paper following
the reviewers' comments.

Below, we outline our responses to the Reviewers' comments and questions. In the revised manuscript, the changes marked in {\color{red}red}
correspond to remarks made by Reviewer \#2, and those marked in {\color{blue}blue} correspond to comments made by Reviewer \#4.

\bigskip
\underline{Response to the Comments of Reviewer \#2}

\begin{itemize}[wide, labelwidth=!, labelindent=0pt]
\item
{\bf Major comments:}

\begin{itemize}[wide, labelwidth=!, labelindent=0pt]
\item{\bf In 3.5, the conservative and nonconservative solutions are blended using parameter s(epsilon). Is blending done after all
calculations are completed or at each time step?}

\smallskip
{\sc\underline{Reply}}: The blending is performed at each stage of the SI-DeC time discretization. We have clarified this point in \S3.5.

\medskip
\item{\bf To see the effect of the present proposal of the dual formulation, two solutions using the conservative and nonconservative
solutions (V-solution and U-solution) can be shown in some of the calculation examples. If so, we will directly see that the interpolation
using s makes a better solution.}

\smallskip
{\sc\underline{Reply}}: We have modified Example 5 to demonstrate the importance on the post-processing for the large (intermediate) values
of $\varepsilon$.
\end{itemize}

\medskip
\item
{\bf Minor comments:}

\begin{itemize}[wide, labelwidth=!, labelindent=0pt]
\item{\bf In p.14, "PCCU" needs full spelling at the first time.}

\smallskip
{\sc\underline{Reply}}: We have fixed it and introduced the abbreviation PCCU on page 2, where it is first mentioned.

\medskip
\item{\bf In 3.4, can Step 1 and 2 be exchanged because the calculation result $\rho_{j,k}^*$ in Step 1 is not used in Step 2.}

\smallskip
{\sc\underline{Reply}}: Yes, indeed. Actually, these two steps can be independently performed in parallel. Our reported algorithm
describes a plain serial implementation, however, some steps can be parallelized. We have added Remark 3.3 on this at the end of \S3.4.

\newpage
\item{\bf Figure 4.1 shows the relation between $s$ and $\varepsilon$ used in the present study. Is it used for all calculation examples?}

\smallskip
{\sc\underline{Reply}}: Yes, we have added Remark 4.2 after the definition of $s(\varepsilon)$ in the beginning of \S4 to prevent future
readers from having the same doubt.

\medskip
\item{\bf In Fig.4.2, the caption of the vertical axis and its definition should be written.}

\smallskip
{\sc\underline{Reply}}: We have revised the manuscript accordingly. In particular, we have added a precise definition of the plotted
quantities and clarified the meaning of the vertical axis in Figure 4.2.

\medskip
\item{\bf Example 4 is not understood well. We can see a large jump of density on $y=4$ at $t=0$. Usually, a large density jump makes wave
propagation in later times but we cannot see such waves at $t=10$ and $20$.}

\smallskip
{\sc\underline{Reply}}: We believe the reviewer is referring to Example 3 (Baroclinic Vorticity Generation), as the density is constant in
Example 4. Indeed, in this test, a visible density jump is present at $y=4$ in the initial condition. This density discontinuity, however,
is not accompanied by a corresponding pressure discontinuity. In this setting, the flow evolution is driven by the interaction of the smooth
pressure and velocity perturbations with the density interface. Since the pressure remains continuous across the interface, the density jump
alone does not generate strong waves. Instead, this interaction generates vorticity and leads to the formation of shear-layer instabilities.

To avoid possible confusion, we have added a short clarification in the revised manuscript.

\medskip
\item{\bf In [2] in the reference list, authors' names are not written.}

\smallskip
{\sc\underline{Reply}}: This is a standard convention of the SIAM bibliography style, where consecutive references by the same author(s) are
abbreviated using a dash. To avoid possible confusion, we have modified the bibliography style so that all authors' names are explicitly
displayed.
\end{itemize}
\end{itemize}

\bigskip
\underline{Response to the Comments of Reviewer \#4}

\begin{itemize}[wide, labelwidth=!, labelindent=0pt]
\item{\bf Evolving two complete sets of variables leads to roughly double the memory requirements and substantially higher computational
cost per time step (two Poisson problems per step).}
	
\smallskip
{\sc\underline{Reply}}: Indeed, there is an overhead with respect to schemes evolving a single set of variables. However, we would like to
point out that this does not strictly correspond to a double computational cost. For example, the same reconstructed cell interface values
are used for both $\bm V$- and $\bm U$-solutions. Moreover, the evolution of the $\bm U$-solution is much computationally cheaper than the
evolution of the $\bm V$-solution, as it is performed with an explicit scheme.

\smallskip
We have added Remark 3.4 on this at the end of \S3.4.
	
\medskip
\item{\bf Additionally, the computation of characteristic wave speeds relies on global quantities $\rho_{\max}$ and $p_{\min}$ (the maximum
density and minimum pressure over the entire domain). In large-scale parallel 3D simulations, this introduces frequent global reductions,
creating a notable communication bottleneck.}
	
\smallskip
{\sc\underline{Reply}}: The computation of the global quantities $\rho_{\max}$ and $p_{\min}$ only involves two scalar reductions and
therefore represents a relatively small overhead compared with the overall cost of the algorithm, which is dominated by the solution of the
linear systems arising from the pressure equation. Nevertheless, we agree that any global communication may affect scalability in extremely
large-scale simulations.
	
\medskip
\item{\bf To prevent wave-speed collapse in higher-Mach regions, the authors introduce an ad-hoc $\epsilon^4$ regularization. While this
modification vanishes in the low-Mach limit, it remains an artificial fix.}

\smallskip
{\sc\underline{Reply}}: We agree that this is an artificial fix, but we found it necessary to improve the handling of discontinuous
solutions in large (intermediate)-Mach-number regimes. Remark 4.1 has been added in the beginning of \S4.

\medskip
\item{\bf The method also employs a heuristic switching function $s(\epsilon)$ with empirically tuned parameters ($\epsilon_0=0.15$,
$\epsilon_1=0.4$, $\alpha=14$) to blend the primitive and conservative solution, across all the numerical illustrations. Room for deeper
mathematical understanding or Machine Learning is ok here ?
}

\smallskip
{\sc\underline{Reply}}: Indeed, the adopted parameters have been selected empirically. Nevertheless, we would like to stress that the wide
range of numerical experiments reported in the manuscript indicates that the proposed choice is sufficiently robust. We agree that a deeper
theoretical understanding of the optimal design of the switching function and its parameters would be highly desirable. 

\smallskip
These aspects are now discussed in Remark 4.2 in the beginning of \S4.
	
\medskip
\item{\bf Another important issue is that the primitive pressure equation (Eq.~2.7) depends on the term $\gamma p\nabla\cdot\mathbf{u}$,
making the entire framework specific to ideal gases. Extending the method to real gases or complex thermodynamics would require a complete
re-derivation of the splitting and the Poisson solver.}
	
\smallskip
{\sc\underline{Reply}}: Indeed, the present derivation is performed for the ideal-gas Euler equations and extending the framework to more
general equations of state would require a re-examination of the primitive-variable formulation including the pressure equation. Extensions
to more complex thermodynamic models are certainly of interest, but lie beyond the scope of the current paper. We have added Remark 3.1 at
the end of \S3.1.1 in the revised manuscript to clarify this point.
	
\medskip
\item{\bf The authors describe purely conservative flux splitting as presenting ``significant challenges.'' While this is understandable,
the claim should be nuanced in light of my previous remarks. The CFD community has already developed several viable and strictly
conservative asymptotic-preserving schemes for the Euler equations.}

\smallskip
{\sc\underline{Reply}}: We fully agree that several viable and fully conservative AP schemes for the Euler equations have been successfully
developed in recent years, and we have mentioned some of the relevant works in the Introduction. We have also emphasized there that the
flux splitting strategy proposed in [26] is, in a certain sense, optimal for the isentropic Euler equations, and that extending this
particular splitting strategy to the full Euler equations is not straightforward.

\smallskip
In this context, the main advantage of the proposed AP DF-FV method is that it can be viewed as an extension of the optimal splitting
strategy from [26] to the full Euler system through the nonconservative formulation.

\medskip
\item{\bf Overall, the manuscript by Chertock et al.\ is a creative and mathematically rigorous contribution. It successfully demonstrates
that dynamically coupling a non-conservative low-Mach solver with a conservative shock-capturing solver can yield a theoretically sound and
asymptotic-preserving method. Nevertheless, the Dual Formulation approach does not completely overcome the practical difficulties of
low-Mach CFD. For large-scale industrial applications, particularly those involving complex thermodynamics or massive parallelism, existing
unified conservative AP methods may offer better efficiency and scalability. The authors may want to consider these remarks in their revised
manuscript.}
	
\smallskip
{\sc\underline{Reply}}: We have added Remarks 3.1, 3.4, and 4.2, in which some of the above points raised by the reviewer have been
mentioned. In fact, we are very interested in extending our DF-FV approach to more practically relevant systems and we plan to investigate
such possibilities in the nearest future.

\medskip
\item{\bf In Definition 3.1, the phrase ``computed solution'' could be misinterpreted as the final discrete solution. For greater clarity,
the authors should explicitly describe $\rho^n$ and $p^n$ as semi-discrete (in time) solutions.}
	
\smallskip
{\sc\underline{Reply}}: We have modified Definition 3.1 accordingly.
\end{itemize}
\vspace*{2cm}
\closing{Sincerely} % eg Regards,
\end{letter}

\end{document}

